%!TeX root=MemoriaTFG.tex

\chapter{Conclusions}
\label{cap:conclusions}

Aquest treball partia d'una motivació doble. D'una banda, un repte tècnic: aplicar
l'\ac{AR} profund a un joc d'informació parcial poc estudiat en la literatura, el Truc, que
hi afegeix apostes encadenades i farol. De l'altra, una motivació personal i cultural:
dotar un joc tradicional amenaçat per la manca de jugadors d'un rival artificial que en
permeti la pràctica encara que no hi hagi prou humans a la taula. Aquests dos fils s'han
desplegat a través dels tres objectius fixats a la introducció: modelitzar i implementar el
joc, comparar algorismes i tècniques d'entrenament, i avaluar el comportament dels agents
resultants. El resultat és un arc que va des d'agents incapaços
de superar el 40\% de win-rate contra una heurística senzilla (Fase~1) fins a un agent
competitiu i estadísticament proper a l'equilibri de Nash (Fase~6).

\section{Principals conclusions}

\paragraph{Modelització i aplicació jugable.} El primer objectiu s'ha assolit plenament.
S'ha construït un simulador fidel del Truc amb arquitectura \ac{MVC}, amb dues interfícies
(consola i escriptori) i un mode humà contra agent, i s'ha exposat el mateix motor com a
entorn estàndard compatible amb les biblioteques d'\ac{AR}, en variants d'un sol agent i
multi-agent. El motor admet fins i tot la modalitat de dos contra dos amb senyes, encara
que l'entrenament d'agents s'hagi circumscrit a la variant d'un contra un. La separació
estricta de responsabilitats ha estat clau: substituir un rival humà per un agent d'\ac{AR}
no ha exigit cap canvi ni al controlador ni al nucli del joc.

\paragraph{Comparació d'algorismes i tècniques.} El segon objectiu deixa una lliçó
transversal: el coll d'ampolla no era el còmput, sinó el \textbf{senyal d'aprenentatge}. El
cas més clar és el \ac{PPO}, l'agent més ràpid i escalable però el més feble sobre partides
senceres; factoritzar la partida en mans individuals (\emph{curriculum learning}) el va
disparar del 3{,}8\% al 84{,}8\% de win-rate ponderat. Les fases següents refinen aquesta
idea amb resultats matisats. Un extractor de característiques convolucional no aporta res
per si sol amb pesos aleatoris, però preentrenat de forma supervisada i congelat arriba al
91{,}2\% (\ac{DQN}): la inducció estructural només compensa amb una inicialització
informativa. La memòria recurrent (\ac{LSTM}) no va rendir dins del pressupost disponible
(82{,}0\% enfront del 86{,}0\% de la variant sense memòria, amb un cost 46 vegades superior),
principalment per un desajust entre l'horitzó d'entrenament i el d'avaluació. Finalment, el
\emph{self-play} mixt inspirat en el \ac{FSP} va millorar la robustesa davant estils
d'oponent diversos (la robustesa inter-estil, \(\text{metric\_robust}\), va pujar de
73{,}2\% a 77{,}8\%) sense degradar el rendiment global, i el \ac{NFSP} complet va assolir
un punt d'equilibri de Nash empíric genuí (\emph{Nash gap} de 0{,}0~pp). Ara bé, el \ac{NFSP}
no va reduir la dispersió entre estils respecte al \emph{self-play}: aproximar l'equilibri de
Nash i maximitzar la robustesa inter-estil no són el mateix eix.

\paragraph{Avaluació del comportament.} El tercer objectiu combina una anàlisi quantitativa
i una de qualitativa. El model desplegat bat totes les variants de l'oponent heurístic
(entre el 55{,}0\% i el 74{,}8\%) i empata contra si mateix (49{,}8\%), amb una política
mixta i poc predictible. Quantitativament, ha après una autèntica divisió del treball entre
les dues apostes: el truc funciona com a farol (desacoblat de la força de la mà) i l'envit
es reserva per a les mans amb punts. Aquesta és, alhora, la conclusió més important del
treball i la seva limitació més interessant: una proximitat estadística a l'equilibri
\textbf{no equival} a un rival creïble i inexplotable per a una persona. El farol al truc,
tot i ser eficaç contra oponents fixos, és tan sistemàtic que resulta llegible, i un jugador
humà atent l'explotaria pagant el truc i esperant. La rúbrica qualitativa ho reflecteix amb
una nota global de 2{,}5 sobre 5: gestió de l'envit sòlida, però farol i adaptació al
marcador febles.

\section{Limitacions}

Els resultats s'han d'interpretar dins d'unes fronteres clares. L'entrenament s'ha limitat a
la variant d'un contra un sense senyes, de manera que la coordinació per parelles (el nucli
social del Truc) queda fora de l'estudi. L'avaluació qualitativa s'ha fet amb un model de
llenguatge com a jutge (\emph{LLM-as-judge}), que és un indicador automàtic i no una
validació amb jugadors humans reals. Algunes fases van treballar amb pressupostos reduïts
(la memòria recurrent, 12M passos enfront dels 24M de les línies base) i amb un únic
oponent de referència o poques llavors, cosa que limita la potència estadística d'alguns
contrastos. Finalment, els valors de rendiment provenen dels experiments propis: no
existeix cap banc de proves extern per al Truc que permeti situar l'agent respecte a
altres sistemes.

\section{Treball futur}

La línia natural de continuació és la variant \textbf{cooperativa de dos contra dos amb
senyes}, ja suportada pel motor però no entrenada. Introdueix dos reptes nous: la
coordinació entre companys sense comunicació verbal explícita i la lectura (i possible
explotació) de les senyes de l'adversari, un escenari genuïnament multi-agent. En segon
lloc, val la pena reprendre la memòria recurrent amb el desajust corregit: alinear
l'horitzó d'entrenament i el d'avaluació (sessions de partides senceres), reduir la mida de
l'\ac{LSTM} i ampliar el pressupost. En tercer lloc, caldria una validació de credibilitat
amb jugadors humans i estratègies específiques per fer el farol al truc menys predictible,
com penalitzar els patrons regulars o adaptar l'agressió al marcador sense introduir
debilitats explotables. Finalment, seria interessant comparar l'enfocament \ac{DRL} amb
mètodes basats en \ac{FSP} tabular o \emph{Counterfactual Regret Minimization} adaptats al
Truc, per contrastar la qualitat de l'equilibri assolit.

En conjunt, el treball demostra que és possible entrenar des de zero un agent competitiu
per a un joc de cartes tradicional d'informació parcial i, sobretot, deixa una lliçó que va
més enllà del Truc: un agent proper a l'equilibri contra oponents heurístics pot continuar
sent llegible per a una persona. Amb tot, l'objectiu que va originar aquest \ac{TFG} queda
acomplert en la variant d'un contra un: el Truc pot continuar viu encara que no hi hagi prou
humans disponibles, i la partida per parelles amb senyes resta com el pas següent per
perpetuar plenament aquesta motivació inicial.
